\subsection{The final coordinate of the UV difference}
\label{sec:comp-uv-final-coordinate}
Throughout this calculation the electron mass is the unit of energy,
$s=P^2>1$, $P=p+k_{\mathrm{in}}$ is the incoming total momentum,
with $k_{\mathrm{in}}$ the incident photon momentum, and $H=s-1$ is
fixed and positive. The original Dalitz variables are
\begin{equation}
 x=2P\cdot k_1,\qquad y=2P\cdot k_2,\qquad
 z=x+y-H=2k_1\cdot k_2.
\end{equation}
The final electron is on shell. Positive final energies and the
three-body kinematic boundary give
\begin{equation}
 0<x<H,\qquad 0<z<\frac{x(H-x)}{s-x},\qquad y=H-x+z.
\end{equation}
The hyperbolic angle in the $UV$ angular function is rationalized by
its positive exponential $v$. Explicitly,
\begin{align}
 D_v&=H^2v-x(v-1)^2,\qquad
 z=\frac{x(H-x)(v-1)^2}{D_v},\\
 y&=\frac{H^2v(H-x)}{D_v},\qquad
 \left|\frac{\partial(x,z)}{\partial(x,v)}\right|
 =\frac{x(H-x)H^2(v^2-1)}{D_v^2}.
\end{align}
At fixed $x$, the lower and upper Dalitz boundaries map to $v=1$ and
$v=s$. The Jacobian is positive in the interior. Thus the new domain
is the rectangle $0<x<H$, $1<v<s$. This is why the inner energy
integration can be completed before the final $v$ integration.
In particular,
\begin{equation}
 xy=\frac{x(H-x)}{1-kx/H},\qquad k=\frac{(v-1)^2}{Hv},
 \qquad 0\leq k\leq\frac Hs<1.
\end{equation}
The analytic regulator therefore has the positive real value
$[x(H-x)]^\eta(1-kx/H)^{-\eta}$. The denominator is at least $1/s$
on the closed rectangle. No logarithmic branch is inferred from a
numerical result.

After the inner energy integration, the $UV$ angular difference sector
is a finite sum of terms
\begin{equation}
 R_{ab}(H,v)\ln^a v\ln^b(1-k),\qquad
 k=\frac{(v-1)^2}{Hv},\qquad 1\leq v\leq H+1.
\end{equation}
Here $H=s-1>0$, and the $R_{ab}$ are the rational functions obtained
from the real-emission amplitude. Their coefficients may include
$\ln H$, which is constant in this integration. The pole and finite
parts of the analytic regulator are treated separately.

\paragraph{The endpoint Laurent coefficients.}
The separated rational energy kernels have poles of at most second
order at $x=0,H$. Other poles are at $x=H/k>H$ and hence outside the
integration path. Start with $\operatorname{Re}\eta>1$, where these
endpoint integrals converge absolutely. Subtract the constant and linear
Taylor terms of the smooth endpoint factor before analytically continuing
to a neighborhood of $\eta=0$. This operation gives an ordinary convergent
remainder and explicit Mellin integrals for the removed terms.
For example, near $x=0$ the smooth factor obeys
\begin{equation}
 (H-x)^\eta(1-kx/H)^{-\eta}
 =H^\eta\left[1+\eta\frac{k-1}{H}x+O(x^2)\right].
\end{equation}
The kernel $x^{-2}$ consequently has the finite coefficient
$-1/H+(k-1)/H=(k-2)/H$. The first term is the finite part of the
unregulated double pole. The second follows from $\eta x$ multiplying
$x^{-2+\eta}$. Although that term vanishes pointwise as $\eta\to0$,
its endpoint integral has a nonzero limit.
At $x=H$, write $b=H-x$. The corresponding expansion is
\begin{equation}
 x^\eta(1-kx/H)^{-\eta}
 =\left(\frac H{1-k}\right)^\eta
 \left[1-\frac{\eta b}{H(1-k)}+O(b^2)\right].
\end{equation}
Including the signs in $(x-H)^{-1}=-b^{-1}$ gives the complete rules
\begin{align}
 \int_0^H\frac{[x(H-x)]^\eta}{x(1-kx/H)^\eta}\,\dd x
 &=\frac1\eta+2\ln H+O(\eta),\\
 \int_0^H\frac{[x(H-x)]^\eta}{x^2(1-kx/H)^\eta}\,\dd x
 &=\frac{k-2}{H}+O(\eta),\\
 \int_0^H\frac{[x(H-x)]^\eta}{(x-H)(1-kx/H)^\eta}\,\dd x
 &=-\frac1\eta+\ln(1-k)-2\ln H+O(\eta),\\
 \int_0^H\frac{[x(H-x)]^\eta}{(x-H)^2(1-kx/H)^\eta}\,\dd x
 &=-\frac{2-k}{H(1-k)}+O(\eta).
\end{align}
The displayed integrals near $\eta=0$ mean their analytic continuations
from the convergence half-plane. The other energy kernels are bounded
on $0\leq x\leq H$ and can be integrated at $\eta=0$ for their finite
coefficient. These rules explain every regulator-generated constant used
in the inner integration. They do not discard power-divergent endpoint
terms by hand.

Put $\zeta=v-1$ for the final polylogarithm coordinate. This $\zeta$ is
distinct from the original Dalitz variable $z$. The two roots of $1-k$ are
\begin{equation}
 r_\pm=\frac{H+2\pm\sqrt{H(H+4)}}2,\qquad
 r_+r_-=1,\qquad 0<r_-<1<H+1<r_+.
\end{equation}
Factorization, together with the value at $\zeta=0$, gives
\begin{align}
 \ln v&=G(-1;\zeta),\\
 \ln(1-k)&=G(r_+-1;\zeta)+G(r_--1;\zeta)-G(-1;\zeta).
\end{align}
The normalization of the multiple polylogarithms is
\begin{align}
 G(\varnothing;\zeta)&=1,\qquad
 G(a,\boldsymbol b;\zeta)
 =\int_0^\zeta\frac{\dd u}{u-a}G(\boldsymbol b;u),\\
 G(\underbrace{0,\ldots,0}_{n};\zeta)&=\frac{\ln^n\zeta}{n!}.
\end{align}
For words with a convergent lower endpoint, the path is the real segment.
The last line fixes the normalization for logarithmic lower endpoints.
All logarithms are real on $0\leq \zeta\leq H$. Their letters lie outside
this interval. Products are expanded by the shuffle identity, and the
rational factors are integrated by partial fractions and the finite
polylogarithmic reduction already described for the soft angular term.

\paragraph{A removable origin.}
Individual rational coefficients can have poles at $\zeta=0$, even though
their sum multiplied by the logarithms is regular. An origin pole must
therefore be handled before setting the lower endpoint to zero. For an
analytic word write
\begin{equation}
 G(\boldsymbol a;\zeta)=\sum_{j\geq1}g_j(\boldsymbol a)\zeta^j.
\end{equation}
Normalize each primitive to have zero finite part at $\zeta=0$, meaning
zero constant term after separating its negative powers and logarithms.
For $\boldsymbol a=(b,\boldsymbol c)$ and $n\geq2$, integration by
parts then gives
\begin{align}
 \mathfrak T[\zeta^{-n};b,\boldsymbol c](\zeta)
 &=-\frac{G(b,\boldsymbol c;\zeta)}{(n-1)\zeta^{n-1}}
 +\frac{g_{n-1}(b,\boldsymbol c)}{n-1}\\
 &\quad+\frac1{n-1}\mathfrak T\!\left[
 \frac1{\zeta^{n-1}(\zeta-b)};\boldsymbol c\right](\zeta).
\end{align}
The displayed constant is essential. It subtracts the finite part of
the boundary term and is absent in the simpler case of a nonzero pole.
For the empty word the origin rules are
\begin{equation}
 \mathfrak T[\zeta^{-1};\varnothing](\zeta)=\ln \zeta,\qquad
 \mathfrak T[\zeta^{-n};\varnothing](\zeta)=\frac{\zeta^{1-n}}{1-n}
 \quad(n\geq2).
\end{equation}
Linearity then makes the normalized primitive of the full regular
integrand vanish at zero. Negative powers and $\ln \zeta$ must cancel in
the complete expression. This is a check on the independently derived
rational coefficients, not a choice of a finite subtraction constant.

An elementary example displays the role of the constant without any
numerical evaluation,
\begin{equation}
 \int_0^Z\frac{\ln(1+\zeta)-\zeta+\zeta^2/2}{\zeta^3}\,\dd\zeta
 =\left(\frac12-\frac1{2Z^2}\right)\ln(1+Z)
 +\frac1{2Z}-\frac14.
\end{equation}
The right-hand side differentiates to the integrand and tends to zero
as $Z\to0^+$. The $-1/4$ is the finite-part normalization constant.

The script \texttt{compton\_soft\_uv\_check.py} implements this reduction.
It reads the two independently generated inner-energy coefficient files,
forms their difference, and evaluates the resulting ordinary multiple
polylogarithms by convergent power series and analytic continuation.
The original coordinate integral is used only as a separate check.
For general $H>0$, direct expansion of the rational coefficients gives
\begin{align}
 \mathcal P_{\mathrm{UV}}(H,1+\zeta)
 &=-\frac{256(H+1)(H+5)}{15H^2}\zeta+O(\zeta^2),\\
 \mathcal F_{\mathrm{UV}}(H,1+\zeta)
 &=-\frac{256(H+1)}{15H^3}
 \bigl[2H(H+5)\ln H-H^2-5H-15\bigr]\zeta
 +O(\zeta^2).
\end{align}
These follow by expanding $\ln(1+\zeta)$ and
$\ln[1-\zeta^2/(H(1+\zeta))]$ before multiplying the rational Laurent
series. The coefficients of $\zeta^{-4},\ldots,\zeta^{-1},\zeta^0$
vanish identically in both sums. Only finitely many Taylor coefficients
are required to prove these cancellations, since the highest rational
origin pole is fourth order.
Thus the apparent origin poles cancel for every $H>0$.
\subsubsection{Constant letters for the UV contribution}
The moving letters of the integrated $UV$ contribution can be removed
without using the total cross section.  Introduce
\begin{equation}
 q=\frac{\sqrt{H+4}-\sqrt H}{\sqrt{H+4}+\sqrt H},\qquad
 H=\frac{(1-q)^2}{q},\qquad 0<q<1.
\end{equation}
Write $z=v-1$.  Its upper endpoint is $H$, and the two logarithms of the
inner Laurent coefficient become
\begin{align}
 \ln v&=G(-1;z),\\
 \ln(1-k)&=G(q-1;z)+G(q^{-1}-1;z)-G(-1;z),
 \qquad k=\frac{z^2}{H(1+z)}.
\end{align}
This follows by factoring $1-k$ and matching its value one at $z=0$.
All factors are positive on $0\leq z\leq H$: the negative letter $q-1$
is outside the interval, and $q^{-1}-1>H$.  Rational partial fractions
and integration by parts therefore produce GPL words with letters
$0,-1,q-1,q^{-1}-1$ and rational coefficients in $q,z$.

To replace their moving endpoint and letters by constants, differentiate
the whole word with respect to $q$.  For a word $w=(a_1,\ldots,a_n)$,
let $w\setminus j$ mean deletion of its $j$th entry.  The total differential
is
\begin{align}
 \dd G(w;H)={}&G(w\setminus1;H)\,\dd\ln(H-a_1)
 -G(w\setminus n;H)\,\dd\ln(-a_n)\\
 &+\sum_{j=1}^{n-1}
 [G(w\setminus(j+1);H)-G(w\setminus j;H)]
                      \,\dd\ln(a_j-a_{j+1}).
\end{align}
The last endpoint term is absent when $a_n=0$.  Adjacent equal letters
are treated by continuity, with their paired difference equal to zero.
Factoring each logarithmic differential gives only constant simple poles.
In this case the complete alphabet is
\begin{equation}
 \left\{0,1,-1,\frac{1+\ii\sqrt3}{2},
                    \frac{1-\ii\sqrt3}{2}\right\}.
\end{equation}
For example, $H+1=(q^2-q+1)/q$ explains the complex-conjugate pair.
Its appearance does not change the real physical boundary value.

The recursive integration uses basepoint $q=1$ approached from below.
For every nonzero moving letter, $H/a_j\to0$ there.  Thus every original
word containing a nonzero letter vanishes, including powers of logarithms
multiplying that vanishing limit.  Pure-zero words instead have
\begin{equation}
 G(\underbrace{0,\ldots,0}_{n};H)=\frac{\ln^nH}{n!},\qquad
 \ln H=2\ln(1-q)-\ln q.
\end{equation}
These relations fix all integration constants.  The recursion checks that
pure powers of $\ln(1-q)$ cancel for every nonzero-letter original word.

The lower endpoint of the angular primitive supplies an additional finite
constant, even though the combined integral is convergent.  With
\begin{align}
 P(q)&=3q^8+18q^7+12q^6-1874q^5+1666q^4
                         -1874q^3+12q^2+18q+3,\\
 Q(q)&=-25q^8+238q^7-436q^6+2274q^5-4294q^4
                         +2274q^3-436q^2+238q-25,
\end{align}
the constants to subtract from the upper primitive are
\begin{align}
 b_{UV,-1}&=-\frac{(q^2-q+1)P(q)}{9q^3(q-1)^4},\\
 b_{UV,0}&=-\frac{q^2-q+1}{18q^3(q-1)^4}
                         [4P(q)\ln H+Q(q)].
\end{align}
The subscripts label the pole and finite Laurent coefficients, respectively.
These constants follow from the Taylor coefficients of each GPL multiplied
by its rational prefactor at $z=0$.  Setting every GPL to zero before
multiplying by the prefactor would lose them.

The resulting pole and finite coefficients contain respectively seven and
41 constant-letter words of weight at most three, before multiplication
by $1/(8sH)$. The script \nolinkurl{compton_soft_uv_constant_basis.py}
generates their explicit coefficient lists,
\nolinkurl{soft_uv_constant_basis_pole.json} and
\nolinkurl{soft_uv_constant_basis_finite.json}, in
\nolinkurl{build/intermediate/}.
The reduction script uses only the independent inner Laurent coefficients.
The primitive follows the exact rational integration recurrence above.
The primitive is checked by differentiating each independent kernel
term before summing. If its coefficient dictionary is $Q_w(y)$,
the derivative follows from
\begin{equation}
 \frac{\mathrm d}{\mathrm dy}\bigl[Q_w(y)G(w;y)\bigr]
 =Q'_w(y)G(w;y)+\frac{Q_w(y)}{y-a}G(\widehat w;y),
 \qquad w=(a,\widehat w).
\end{equation}
For the empty word only the first term occurs. Subtracting the input
kernel leaves a rational coefficient for each GPL word.
Every coefficient vanishes exactly, for both pole and finite parts.
The calculation holds the external parameter $q$ and $\ln H$ fixed
while differentiating in $y=v-1$.
Canceling rational factors before differentiation makes this test finite
in size. Linearity then proves the derivative identity for their sum.
The reproducible checks are in
\texttt{compton\_uv\_primitive\_certificate.py}.
Evaluation against the earlier moving-letter expression provides
a separate numerical check of the conversion.


\subsection{The UW endpoint matching}
Use the original Dalitz coordinates defined above, and set
$a=H-x$. The positive angular variable satisfies $1<v<1+a$, so
$t=(v-1)/a$ places its two endpoints at zero and one. The complete map is
\begin{align}
 x&=H-a,\qquad v=1+at,\qquad
 z=\frac{a(H-a)t^2}{1+at},\\
 y&=\frac{a[1+Ht^2+at(1-t)]}{1+at},\\
 \left|\frac{\partial(x,z)}{\partial(a,t)}\right|
 &=\frac{a(H-a)t(2+at)}{(1+at)^2}>0.
\end{align}
The new domain is exactly $0<a<H$, $0<t<1$. Both the sign of the
Jacobian and the range of $t$ follow from the positive solution for $v$.
In these rectangular coordinates the analytic
regulator is $[a(H-a)f(a,t)]^\eta$, with
\begin{equation}
 f(a,t)=\frac{1+Ht^2+at(1-t)}{1+at},\qquad
 0<a<H,\quad 0<t<1.
\end{equation}
The useful alternative form $f=1+(H-a)t^2/(1+at)$ shows that
$1\leq f\leq1+H$ on the rectangle. Its logarithm and endpoint derivatives
are real and bounded for fixed $H>0$. For a fixed interior $t$, start
again at $\operatorname{Re}\eta>1$ and continue analytically only after
subtracting the endpoint Taylor terms.
Write the unregulated difference as
$F(a,t)=L(a,t)\ln(1+at)-T(a,t)$, where $L$ and $T$ are the rational
coefficient and Taylor subtraction obtained after including the Jacobian.
At the upper endpoint put $b=H-a$ and expand
\begin{equation}
 L=\frac{L_2}{b^2}+\frac{L_1}{b}+O(1),\qquad
 T=\frac{T_2}{b^2}+\frac{T_1}{b}+O(1).
\end{equation}
These coefficients are fixed algebraically, for example
$L_2=[b^2L]_{b=0}$ and $L_1=[\partial_b(b^2L)]_{b=0}$.
Because $\ln[1+(H-b)t]=\ell-bt/(1+Ht)+O(b^2)$, the complete singular
coefficients are
\begin{align}
 g_2&=L_2\ell-T_2,\qquad \ell=\ln(1+Ht),\\
 g_1&=L_1\ell-\frac{L_2t}{1+Ht}-T_1.
\end{align}
In particular, the derivative of the logarithm contributes to the simple
pole and cannot be omitted. Thus $F=g_2/b^2+g_1/b+O(1)$.
Since $f(H,t)=1$, direct expansion gives
\begin{equation}
 a^\eta f(a,t)^\eta
 =H^\eta\left[1+\eta b\left(-\frac1H+\frac{t^2}{1+Ht}\right)
 +O(b^2)\right].
\end{equation}
The ordinary finite part used for comparison is defined explicitly by
\begin{equation}
 \operatorname{FP}\int_0^H F(a,t)\,\dd a
 =\lim_{\delta\to0^+}\left[
 \int_0^{H-\delta}F(a,t)\,\dd a-\frac{g_2}{\delta}
 +g_1\ln\delta\right].
\end{equation}
After the two singular terms are subtracted, the remainder is bounded
at $b=0$. Its integral is analytic in $\eta$ near zero. The removed terms
are integrated before expansion using
\begin{equation}
 \int_0^H b^{\eta-2}\,\dd b=\frac{H^{\eta-1}}{\eta-1},
 \qquad
 \int_0^H b^{\eta-1}\,\dd b=\frac{H^\eta}{\eta},
\end{equation}
first in their convergence half-planes and then by continuation.
The double-pole kernel therefore does not by itself produce a pole at
$\eta=0$. Its term proportional to $\eta b$, however, multiplies the
second integral and supplies a finite constant. The simple pole produces
$g_1/\eta$. Relative to the ordinary finite
part of the $a$ integral, the finite matching contribution is
\begin{equation}
 g_1\ln H+g_2\left(-\frac1H+\frac{t^2}{1+Ht}\right).
\end{equation}
The second term arises because the factor $\eta b$ multiplying
$b^{-2+\eta}$ leaves a finite limit after integration. It must not be
lost by setting $\eta=0$ before integrating. The finite part of the
unregulated integral already includes $-g_2/H$ and $g_1\ln H$.
The matching term is additional to those terms. At $a=0$ the possible
simple-pole coefficient cancels in the complete angular difference.
Explicitly, $\ln(1+at)=at+O(a^2)$ implies
\begin{equation}
 \lim_{a\to0}aF(a,t)
 =t\lim_{a\to0}a^2L(a,t)-\lim_{a\to0}aT(a,t)=0.
\end{equation}
The last equality is an exact identity of the rational coefficients.
Therefore there is no matching term from that endpoint.

The cancellation at the other boundary can also be seen before numerical
evaluation. Expanding $\ln(1+Ht)$ and the explicit rational endpoint
coefficients gives, for every fixed $H>0$,
\begin{align}
 g_2(H,t)&=\frac{64H^2(H+1)}{15}t+O(t^2),\\
 g_1(H,t)&=-\frac{64H(H+1)}5t+O(t^2).
\end{align}
All the coefficients of $t^{-4},t^{-3},t^{-2},t^{-1},t^0$ vanish in
these complete combinations. Thus the functions entering the pole and
finite endpoint matching are integrable at the lower angular boundary.
This algebraic cancellation, rather than their numerical values at a
few points, justifies treating that boundary as removable.

The remaining $t$ integral of $g_1$ is rational times $\ln(1+Ht)$.
The same finite-part polylogarithmic reduction gives seven terms.
\paragraph{Logarithmic lower endpoints.}
For the finite coefficient, the inner integral also contains $\ln t$
and weight-two functions. A word ending in zero can be rewritten by
shuffle identities. In particular,
\begin{equation}
 G(a,0;t)=\ln t\,G(a;t)-G(0,a;t).
\end{equation}
For a longer word the same identity recursively removes its trailing
zeros. The result is a polynomial in $\ln t$ whose coefficients are
regular-origin polylogarithms. In the repeated-pole primitive, the
finite-part constant is the coefficient of $t^{n-1}(\ln t)^0$.
Terms containing $\ln t$ are therefore included in the reduction rather
than incorrectly treated as ordinary Taylor coefficients.
For example, the shuffle identity gives the exact check
\begin{equation}
 \int_0^1\frac{\ln t}{1+t}\,\dd t
 =G(-1,0;1)=-G(0,-1;1)=\operatorname{Li}_2(-1)
 =-\frac{\pi^2}{12}.
\end{equation}
Differentiation proves the primitive, and its lower endpoint vanishes.

For the finite inner result, use the following names for its eight
logarithmic and dilogarithmic functions:
\begin{align}
 \ell&=\ln(1+Ht),\qquad l_H=\ln H,\qquad
 l_2=\ln(2+Ht),\\
 l_P&=\ln(1+Ht-Ht^2),\qquad l_t=\ln t,\\
 d_1&=\operatorname{Li}_2(-Ht),\qquad
 d_2=\operatorname{Li}_2(-1-Ht)+\frac{\pi^2}{12},\\
 d_3&=\operatorname{Li}_2\!\left(-\frac{(1-t)(1+Ht)}t\right)
       -\operatorname{Li}_2\!\left(-\frac{1-t}t\right).
\end{align}
All logarithm arguments are positive for $H>0$ and $0<t<1$.
The dilogarithm arguments are nonpositive, so their real branches
follow from the original physical interval. Expanding the boundary
formulas in \eqref{eq:comp-real-UW-inner-closed} in these functions
gives thirteen monomials. Write $R_f(H,t)$ for the rational coefficient
of a monomial $f$, including $R_1$ for the constant monomial.

In the complete $UW$ finite sum, the coefficient of
$\ell\ln t$, with $\ell=\ln(1+Ht)$, is exactly the negative of the
coefficient of the dilogarithm difference $d_3$.
This is an identity of the rational coefficients for arbitrary $H$ and
$t$, not a numerical cancellation. It groups these terms into
$d_3-\ell\ln t$. Writing $a=-1/H$ and
$t_\pm=(1\pm\sqrt{1+4/H})/2$, the group is
\begin{equation}
 d_3-\ell\ln t
 =-\sum_{r=t_\pm}\left[G(1,r;t)+G(a,r;t)-G(0,r;t)\right]
 -G(0,a;t).
\end{equation}
Every word on the right has a regular lower endpoint. This identity
also explains why the apparent $\ln t$ terms disappear from the
complete lower-endpoint expansion. The general finite-part rules remain
useful when these terms are integrated separately.

\paragraph{The finite UW coordinate integral.}
Write the inner finite result as a finite sum
$\sum_j R_j(t)M_j(t)$, where the $R_j$ are its explicit rational
coefficients and each $M_j$ is a monomial in the eight logarithmic or
dilogarithmic functions. Translating these functions to ordinary
polylogarithms and applying shuffle products gives finite coefficients
$c_{j\boldsymbol a}$ such that
\begin{equation}
 M_j(t)=\sum_{\boldsymbol a}c_{j\boldsymbol a}
 G(\boldsymbol a;t).
\end{equation}
The final coordinate integral is therefore the finite endpoint expression
\begin{equation}
 \mathcal F_{UW}(H)=\lim_{Z\to1^-}
 \sum_{j,\boldsymbol a}c_{j\boldsymbol a}
 \mathfrak T[R_j;\boldsymbol a](Z).
 \label{eq:comp-real-UW-final-primitive}
\end{equation}
All rational integrations in this formula are eliminated by the terminating
partial-fraction rules, including the origin constants described above.
The functions with a first letter equal to one are combined before the
endpoint limit. Their separate logarithmic divergences need not cancel
within a single term.

For evaluation, the complete inner finite function is analytic at $t=1$.
Indeed, its logarithm arguments are nonzero there and the two arguments
in the dilogarithm difference both tend to zero. The apparent rational
poles cancel in the summed expression. With $d=1-t$, the remaining tail
is consequently obtained algebraically,
\begin{equation}
 \mathcal F_{UW}^{\mathrm{inner}}(H,1-d)=\sum_{n\geq0}c_n d^n,
 \qquad
 \int_{1-\delta}^1\mathcal F_{UW}^{\mathrm{inner}}(H,t)\,\dd t
 =\sum_{n\geq0}\frac{c_n\delta^{n+1}}{n+1}.
\end{equation}
The coefficients are generated by rational power-series division and the
Taylor series of the displayed logarithms and dilogarithms. No phase-space
quadrature enters this endpoint evaluation. Individual terms near $t=0$ require powers of $t$ multiplied by
powers of $\ln t$. The combined function is analytic there, as follows
from the cancellation of $\ln t$ established above.
The generated finite expression requires 52 nonempty polylogarithm words.
\paragraph{The order of the limiting operations.}
All endpoint statements here concern a fixed $H>0$. They do not justify
interchanging the photon-mass limit or the regulator expansion with the
threshold limit $H\to0$. The analytic regulator is first introduced in
a common convergence half-plane. Only after the endpoint Taylor terms
have been isolated are their Mellin integrals continued to $\eta=0$.
The bounded remainders can then be expanded under the integral.

There is also a uniform reason for the removable angular origins. The
hyperbolic function entering each angular integral is analytic at
$q=0$, where
\begin{equation}
 q=\cosh(\ln v)-1=\frac{(v-1)^2}{2v}.
\end{equation}
Together with the Taylor expansion of its rational prefactors, its
constant and linear terms give precisely the terms removed in $J-T$. For the $UV$ map, $q=\zeta^2/[2(1+\zeta)]$ is independent
of the inner energy coordinate. For the $UW$ map,
$q=a^2t^2/[2(1+at)]$ is uniformly small as $t\to0$ for $0\leq a\leq H$.
Taylor's remainder is therefore uniform after the explicit soft endpoint
powers have been factored out. Those powers are at most second order,
and are controlled by $\operatorname{Re}\eta>1$ before continuation.
The local Mellin subtractions above preserve this uniform angular
behavior because their coefficient functions and the regulator shape
factors are analytic in the same angular neighborhood.
This gives the reason for taking the Laurent coefficients before the
last angular integration. A pointwise expansion of the un-subtracted
endpoint integrand would not give that justification.

At the upper $UW$ boundary, the logarithm arguments are positive and
the two $d_3$ dilogarithm arguments tend to zero. At nonzero electron mass, the original hard angular denominators have
no pole on the integration sphere for positive photon energies. The
soft endpoint factors have already been isolated explicitly. Thus the
complete soft-subtracted energy integral has an analytic neighborhood
of $t=1$, even when a particular angular partial-fraction representation
becomes singular there.
The largest apparent rational pole in the separated coefficient list is
fifth order. To determine the tail through order $d^n$, the atom series
are therefore expanded through at least order $d^{n+5}$ before summing.
All negative powers are canceled before the endpoint value is used.
The resulting convergent power series may be integrated term by term
inside its disk of convergence. These arguments establish the endpoint
operations. The numerical verification chapter
tests their implementation independently of these analytic arguments.

\subsubsection{Exact removal of the finite UW endpoint poles}
\label{sec:comp-UW-upper-certificate}
The cancellation of the apparent poles at $t=1$ can be checked for
arbitrary $H>0$, without evaluating the remaining integral numerically.
Write $d=1-t$, $E=\ln(1+H)$, and $C=(H+1)/H$ in this subsection.
Here $E$ is a logarithm, not an electron energy.
Denote the rational coefficient multiplying an atom $f$ in the
thirteen-term inner result by $R_f$.  Only six terms can have a pole
at this endpoint:
\begin{equation}
 R_1+R_{d_3}d_3+R_{l_P}l_P
       +R_\ell\ell+R_{\ell l_t}\ell l_t+R_{\ell l_P}\ell l_P.
\end{equation}
Their rational pole orders are respectively $4,3,5,2,3,3$.
The remaining seven terms have rational denominators nonzero at $t=1$
and analytic atoms there.  Thus they cannot contribute a negative power
of $d$.

The required Taylor coefficients follow from
\begin{align}
 \ell&=E-\sum_{n=1}^4\frac1n
                   \left(\frac{H}{1+H}\right)^n d^n+O(d^5),\\
 l_t&=-\sum_{n=1}^4\frac{d^n}{n}+O(d^5),\\
 l_P&=\sum_{n=1}^4\frac{(-1)^{n+1}}n
                         (Hd-Hd^2)^n+O(d^5),\\
 d_3&=\sum_{n=1}^4\frac{z_1^n-z_2^n}{n^2}+O(d^5),
 \qquad z_1=-\frac{d(1+H-Hd)}{1-d},\quad z_2=-\frac d{1-d}.
\end{align}
Terms of degree above four in the last two finite sums are simply omitted
when constructing the Taylor polynomial.  The two dilogarithm arguments
vanish at $d=0$, so their ordinary convergent series applies.  All these
statements are at fixed $H>0$ in a sufficiently small neighborhood of
$d=0$.  There is no continuation through a dilogarithm cut.

For completeness, the rational expansion requires only polynomial
arithmetic.  If $R_f(1-d)=(-d)^{-m}N(d)/D(d)$ with $D(0)\ne0$, write
$N=\sum N_jd^j$, $D=\sum D_jd^j$.  Its regular quotient coefficients obey
\begin{equation}
 r_0=\frac{N_0}{D_0},\qquad
 r_n=\frac{N_n-\sum_{j=1}^nD_jr_{n-j}}{D_0}.
\end{equation}
Multiplying this finite quotient series by the atom series determines
every negative power.  The largest pole is fifth order, so degree four
is sufficient.  No infinite-series truncation error enters a coefficient
of a negative power.

Let $\operatorname{PP}_{d=0}$ denote the sum of strictly negative powers
in a Laurent expansion.  Substituting the actual rational coefficients
of the inner result gives the particularly simple pair
\begin{align}
 \operatorname{PP}_{d=0}(R_1+R_{l_P}l_P)
 &=C\left[-\frac{16}{d^2}+\frac{4(H-10)}d\right],\\
 \operatorname{PP}_{d=0}(R_{d_3}d_3)
 &=C\left[\frac{16}{d^2}-\frac{4(H-10)}d\right].
\end{align}
The apparent fourth- and third-order poles have already canceled within
the first line.  The fifth-order coefficient vanishes because $l_P(0)=0$.
For the other three terms, divide their principal parts by $8C$.
The coefficients of $d^{-2}$ and $d^{-1}$ are
\begin{center}
\begin{tabular}{lll}
\toprule
Term & $[d^{-2}]/(8C)$ & $[d^{-1}]/(8C)$\\
\midrule
$R_\ell\ell$ & $2E(H+1)/H$ & $-E(H-4-7/H)-2$\\
$R_{\ell l_t}\ell l_t$ & $-2E/H$ & $-7E/H+2/(H+1)$\\
$R_{\ell l_P}\ell l_P$ & $-2E$ & $E(H-4)+2H/(H+1)$\\
\bottomrule
\end{tabular}
\end{center}
Each column sums to zero as a rational identity in $H,E$.  Their
third-order coefficients vanish because $l_t(0)=l_P(0)=0$.
Combining the two cancellations proves
\begin{equation}
 [d^{-j}]F_{UW}(H,1-d)=0,\qquad j=1,\ldots,5,\qquad H>0.
\end{equation}
Here $F_{UW}(H,t)$ denotes the finite inner integrand before the final
$t$ integration.  Since each atom is analytic at the endpoint and the
rational factors have finite pole order, the disappearance of their
complete principal part proves that this endpoint is removable.
It is an identity for the finite integrand, not only for its infrared pole.

The independent script
\nolinkurl{compton_soft_uv_uw_endpoint_certificate.py} reads the thirteen
rational coefficients, computes the six contributing expansions, and
checks every coefficient of $E$ separately using exact arithmetic.
\nolinkurl{uw_endpoint_certificate.json} records the individual terms
as well as the zero sums.  This certificate establishes endpoint
regularity.  It does not provide a bound on the entire omitted numerical
Taylor tail used when evaluating the final special functions.

\subsubsection{Exact removal of the finite UW origin poles}
\label{sec:comp-UW-lower-certificate}
The lower endpoint also admits an exact certificate for every fixed
$H>0$.  The only explicit $\ln t$ in the thirteen-term expression comes
from $d_3$ and the product $\ell\ln t$.  Their rational coefficients
satisfy
\begin{equation}
 R_{d_3}+R_{\ell l_t}=0.
\end{equation}
It is therefore essential to combine these two terms first.  Define
\begin{equation}
 \mathcal D(t)=d_3(t)-\ell(t)\ln t,\qquad \ell(t)=\ln(1+Ht).
\end{equation}
The logarithm cancels exactly, rather than merely becoming integrable.

Its boundary constant follows from the real inversion identity
\begin{equation}
 \operatorname{Li}_2(-z)=-\frac{\pi^2}{6}-\frac12\ln^2z
                              -\operatorname{Li}_2(-1/z),\qquad z>0.
\end{equation}
Apply it to $z=(1-t)(1+Ht)/t$ and $z=(1-t)/t$.  Both are positive on
$0<t<1$, so the same negative-real-axis boundary value is used in both
terms.  The exact result is
\begin{align}
 \mathcal D(t)={}&-\ell\ln(1-t)-\frac12\ell^2
 -\operatorname{Li}_2\!\left(-\frac{t}{(1-t)(1+Ht)}\right)
 +\operatorname{Li}_2\!\left(-\frac{t}{1-t}\right).
\end{align}
Every term on this line is analytic at $t=0$ and vanishes there.
This proves $\mathcal D(0)=0$ independently of any numerical evaluation.
Differentiating either this expression or the original pair gives
\begin{equation}
 \mathcal D'(t)=\frac{l_P}{1-t}-\frac{Hl_P}{1+Ht}
                         +\frac{l_P-\ell}{t},\qquad
 l_P=\ln(1+Ht-Ht^2).
\end{equation}
The apparent $1/t$ is removable because $l_P(0)=\ell(0)=0$.
The resulting series starts as
\begin{equation}
 \mathcal D(t)=-\frac{H^2}{2}t^2
 +\frac{H^2(H+1)}2t^3
 -\left(\frac{11H^4}{24}+\frac{2H^3}{3}\right)t^4+O(t^5).
\end{equation}
In particular, $d_3=Ht\ln t+O(t^2[1+|\ln t|])$ has a zero constant,
fixed by the inversion identity, not chosen to enforce cancellation.

After this combination the entire kernel contains only atoms analytic
at the origin.  Its rational factors have poles of order at most five.
Using the finite quotient recurrence of the preceding subsection,
expand the atom series through degree six.  This determines all negative
powers, the constant, and the coefficient of $t$ exactly.
The cancellation can be displayed without a long list of individual
rational coefficients.  Define the three groups
\begin{align}
 \mathcal B_0&=R_1+R_{l_2}l_2+R_{l_P}l_P+R_\ell\ell,\\
 \mathcal B_1&=R_{d_3}\mathcal D+R_{\ell l_P}\ell l_P+R_{\ell^2}\ell^2,\\
 \mathcal B_H&=\ln H\,[R_{l_H}+R_{\ell l_H}\ell].
\end{align}
Their expansions obtained from the explicit coefficient table are
\begin{align}
 \mathcal B_0={}&(H+1)\left[-\frac{128}{H^2t^2}+
 \frac{64}{Ht}+\frac{4(3H^2-67H-72)}{3H}\right]+O(t),\\
 \mathcal B_1={}&(H+1)\left[\frac{128}{H^2t^2}-
 \frac{64}{Ht}-\frac{4(3H^2-67H-72)}{3H}\right]+O(t),\\
 \mathcal B_H={}&O(t).
\end{align}
Their fifth-, fourth-, and third-order poles have already canceled
inside these groups.  The remaining three atom terms are individually
$O(t)$ because their rational coefficients are regular and their atoms
vanish.  Thus all negative powers and the constant cancel in the full
finite kernel.  The first surviving coefficient is
\begin{equation}
 F_{UW}(H,t)=-\frac{32H(H+1)}{15}(12\ln H-1)t+O(t^2),
                    \qquad t\to0^+,\quad H>0\text{ fixed}.
\end{equation}
There is no residual $\ln t$ in this combined expression.  The origin is
an ordinary analytic endpoint, so the finite lower boundary of the GPL
primitive equals the lower limit of the convergent integral.

The script \nolinkurl{compton_soft_uv_uw_origin_certificate.py} first
checks the exact rational coefficient identity removing $\ln t$.
It then treats $\ln H$ and $\ln2$ as independent symbols and checks
the coefficients of $t^{-5},\ldots,t^{-1},t^0$ separately in exact
arithmetic.  Their values and the nonzero coefficient of $t$ are stored
in \nolinkurl{uw_origin_certificate.json}.  Together with the upper
certificate, this establishes both finite endpoints for arbitrary
$H>0$.  A rigorous error bound for the numerical special-function
series is a separate issue.

\subsection{An exact certificate for the integrated pole sum}
This additional algebraic check tests the separately implemented sector
primitives for every $H>0$.  The physical soft factorization was established
earlier.  No numerical comparison or published total enters this check.
The $UV$ pole density already uses $1<v<H+1$.  In the $UW$ pole density,
set $v=1+Ht$, including $\dd t=\dd v/H$.  The physical infrared angular
integral has the same range after the substitution
\begin{align}
 T&=-t_{\mathrm{elastic}}=\frac{(v-1)^2}{v},
 &y_B&=H-T,\\
 b_e&=\frac{v^2+1}{v^2-1}\ln v,
 &\left|\frac{\dd c}{\dd v}\right|
 &=\frac{2(H+1)}{H^2}\left(1-\frac1{v^2}\right).
\end{align}
Here $y_B=1-u_{\mathrm{elastic}}$ distinguishes the elastic invariant
from the real-emission Dalitz variable. Let
$\mathcal B(v)=\overline{|\mathcal M_0|^2}/e_{\mathrm T}^4$ be the
Born square with its coupling factor removed, evaluated at $x=H$ and
$y=y_B$. The target density for $8H\mathcal C_{\mathrm{IR}}$
is consequently
\begin{equation}
 P(v)=\frac{16(H+1)}H\mathcal B(v)
 \left[\frac{v^2+1}{v^2}\ln v-1+\frac1{v^2}\right].
\end{equation}
Subtract this density from the sum of the $UV$ density and twice the
transformed $UW$ density.  The result is $L(v)\ln v+Q(v)$, with rational
$L$ and $Q$ fixed directly by the sector coefficient tables.  Integration by parts separates rational total derivatives from
simple-pole remainders. This operation, called rational Hermite reduction,
provides a short certificate for its integral.
Define
\begin{align}
 D(v)&=v^2-(H+2)v+1,\\
 p&=H^3-5H^2-28H-24,\qquad d=H^2+6H+4,\\
 j_0&=H^5+4H^4-12H^3+48H^2-32H-384,\\
 j_1&=2H^4+16H^3-40H^2+64H-192,\\
 j_2&=3H^4-16H^3-184H^2-736H-480.
\end{align}
Two rational functions suffice:
\begin{align}
 \mathcal R(v)&=\frac{H+1}{H^3}\left\{
 \frac{128H^2}{(v-1)^2}+16Hv^2
 +\frac{128H(H+2)(v-1)}{D(v)}
 +\frac{32H(3H-4)}{v-1}\right.\\
 &\hspace{14mm}\left.
 -\frac{16H}{v^2}-32dv
 +\frac{192(H+2)^2}{(v+1)^2}
 -\frac{128(H+2)^2}{(v+1)^3}
 +\frac{32p}{v+1}+\frac{32d}{v}\right\},\\
 \mathcal J(v)&=\frac{H+1}{3H^3}\left\{
 \frac{H^3(H+2)(H+6)(H^2-4H+8)[(H+2)v-1]}{D(v)^2}
 +\frac{H^2(j_0+j_1v)}{D(v)}\right.\\
 &\hspace{12mm}
 -\frac{384H^2}{v-1}-72Hv^2-\frac{72H}{v^2}
 -(6H^4-368H^2-1472H-960)v\\
 &\hspace{12mm}\left.
 -\frac{1152(H+2)^2}{v+1}
 +\frac{1152(H+2)^2}{(v+1)^2}-\frac{2j_2}{v}\right\}.
\end{align}
The four constant coefficients needed below are
\begin{align}
 C&=\frac{32(H+1)(H^2-4H-6)}{H^2},\\
 c_1&=\frac{2(H+1)}{3H^3}
 (3H^5-24H^4-112H^3+608H^2+2368H+1920),\\
 c_+&=\frac{256(H+1)p}{3H^3},\qquad
 c_D=-\frac{2(H+1)(H^3-8H^2-8H+48)}H.
\end{align}
Direct differentiation and rational cancellation give the two identities
\begin{align}
 L&=\mathcal R'+\frac Cv,\\
 Q-\frac{\mathcal R}{v}
 &=\mathcal J'+\frac{c_1}{v}+\frac{c_+}{v+1}
                      +c_D\frac{D'}D.
\end{align}
Thus an elementary primitive of the full difference is
\begin{equation}
 \Phi(v)=\mathcal R(v)\ln v+\frac C2\ln^2v+\mathcal J(v)
             +c_1\ln v+c_+\ln(v+1)+c_D\ln[-D(v)].
\end{equation}
There is no unresolved branch in this expression.  On $1\leq v\leq H+1$
the convex quadratic $D$ is negative, since its endpoint values both
 equal $-H$.  All displayed logarithm arguments are positive.
In particular, the $c_D$ logarithm has identical endpoint values.

The apparent pole at $v=1$ must be combined before taking its limit.
If the coefficients of $(v-1)^{-2}$ and $(v-1)^{-1}$ in $\mathcal R$
are $r_{-2}$ and $r_{-1}$, the constant from
$\mathcal R\ln v$ is $r_{-1}-r_{-2}/2$.  Adding the constant from
$\mathcal J$ gives
\begin{align}
 j_-&=-\frac{H+1}{3H^3}
       (H^5+24H^4-28H^3-624H^2-1936H-768),\\
 j_+&=\frac{1}{3H^3(H+1)}
 \left(H^9+6H^8-18H^7-36H^6+351H^5\right.\\
 &\hspace{35mm}\left.+2912H^4+6572H^3+6800H^2+3472H+768\right),\\
 r_+&=-\frac{16(H^5+20H^4+78H^3+114H^2+72H+16)}{H^3(H+1)}.
\end{align}
Here $j_+=\mathcal J(H+1)$ and $r_+=\mathcal R(H+1)$.
With $L_s=\ln(H+1)$ and $L_2=\ln[(H+2)/2]$, the integrated difference is
\begin{equation}
 E(H)=j_+-j_-+(r_++c_1)L_s+\frac C2L_s^2+c_+L_2.
\end{equation}
The remaining rational, single-logarithm, and Bose-restoration pole
expressions contain only the additional logarithms
\begin{equation}
 \ln\frac1{H+1}=-L_s,\qquad
 \ln\left(1+\frac{H^2}{4(H+1)}\right)=2L_2-L_s.
\end{equation}
Substitution and polynomial collection in the two independent symbols
$L_s,L_2$ gives zero for every rational coefficient in
$E+A_{\mathrm{rat}}+A_{\log}+A_D$.  Therefore
\begin{equation}
 A_{\mathrm{rat}}+A_{\log}+A_{UV}+2A_{UW}+A_D
              =8H\mathcal C_{\mathrm{IR}},\qquad H>0.
\end{equation}
The script \nolinkurl{compton_soft_uv_pole_certificate.py} verifies both
rational derivative identities, the endpoint constants, and the final
zero polynomial using exact arithmetic.  It also verifies that the
combined $UV$ and $UW$ pole densities start at order $v-1$, so their
lower endpoints are ordinary convergent integrals.  The machine-readable
certificate is \nolinkurl{soft_pole_certificate.json} in this supplement.
